\documentclass[11pt]{article}
\usepackage[letterpaper,margin=1in]{geometry}
\usepackage[T1]{fontenc}
\usepackage{amsmath,amssymb,amsthm,lmodern,microtype}
\usepackage[hidelinks]{hyperref}

\newtheorem{theorem}{Theorem}
\newtheorem{lemma}[theorem]{Lemma}
\newtheorem{proposition}[theorem]{Proposition}

\title{Tuza's Inequality for Triangle Packing Number at Most Seven}
\author{}
\date{}

\begin{document}
\maketitle

\begin{abstract}
Let $\nu(G)$ denote the maximum number of edge-disjoint triangles in a finite simple graph $G$, and let $\tau(G)$ denote the minimum number of edges meeting every triangle of $G$. We prove that every finite simple graph $G$ with $\nu(G)\leq 7$ satisfies Tuza's inequality $\tau(G)\leq 2\nu(G)$. Building on Kandasamy's maximum-degree-ten theorem, Puleo's reducible configurations, local link reductions, and the Tutte--Berge and Gallai--Edmonds matching theorems, the proof resolves all graph orders.
\end{abstract}

\begin{theorem}\label{thm:main}
Let $G$ be a finite simple graph. If $\nu(G)\leq 7$, then $\tau(G)\leq 2\nu(G)$.
\end{theorem}

The argument uses Kandasamy's maximum-degree-ten theorem~\cite{Kandasamy}, Puleo's reducible sets~\cite{Puleo}, the Tutte--Berge and Gallai--Edmonds matching theorems~\cite{Schrijver}, and the companion result for $\nu(G)\leq 6$~\cite{PackingSix}.

\section{Local reductions and packings at isolates}

Suppose Theorem~\ref{thm:main} fails, and choose a counterexample $G$ minimizing $|E(G)|$ and, subject to that, minimizing $|V(G)|$. Every proper subgraph of $G$ satisfies $\tau\leq 2\nu$. By the packing-six theorem~\cite{PackingSix}, $\nu(G)=7$, so exhibiting either eight edge-disjoint triangles or a triangle cover of size at most fourteen yields a contradiction.

\begin{lemma}\label{lem:basic}
The graph $G$ has no isolated vertices, every edge lies in at least two triangles, and every link $L_v=G[N(v)]$ satisfies $\delta(L_v)\geq 2$. Moreover, for any $v\in V(G)$, any $W\subseteq N(v)$, any nonempty matching $M$ inside $L_v[W]$, and any $Q\subseteq N(v)$ meeting every edge of $L_v$ incident with $W$ outside $M$, we have $|Q|>|M|$.
\end{lemma}

\begin{proof}
An isolated vertex or an edge in no triangle could be deleted without changing $\nu$ or $\tau$. If an edge lies in a single triangle $T$, removing $E(T)$ reduces $\nu$ by at least one and $\tau$ by at most two (the other two edges of $T$ cover all lost triangles), contradicting minimality. Hence every edge lies in at least two triangles, and $\delta(L_v)\geq 2$.

For the scoped link reduction, suppose $|Q|\leq |M|$. Deleting the spokes from $v$ to $W$ and the edges of $M$ yields a proper subgraph $H$. The $|M|$ triangles $\{v\}\cup e$ ($e\in M$) extend any packing of $H$, while $M$ and the spokes from $v$ to $Q$ cover every triangle of $G$ missing from $H$ at cost at most $2|M|$, contradicting minimality.
\end{proof}

Taking $W=N(v)$ in Lemma~\ref{lem:basic} shows that no nonempty matching $M$ in $L_v$ leaves $L_v-M$ with vertex cover number at most $|M|$.

\begin{lemma}\label{lem:puleo}
We have $\delta(G)\geq 5$, the vertices of degree at most six in $G$ form an independent set, each of their links has complement of maximum degree at most one and edge count different from two, and any vertex subsuming a degree-five vertex (meaning its closed neighborhood contains the degree-five vertex's closed neighborhood) has degree at least nine in $G$. Furthermore, every link of $G$ has at least eight edges.
\end{lemma}

\begin{proof}
A link component of order three is a triangle and reduces by Lemma~\ref{lem:basic} using one edge and its opposite vertex. A link component of order four has a two-edge matching and reduces using that matching and the endpoints of one of its edges. Thus $G$ is robust in Puleo's sense, with every link component of order at least five. Since every proper subgraph of $G$ satisfies $\tau\leq 2\nu$, the graph $G$ contains no reducible set in the sense of Puleo's Definition~2.1~\cite{Puleo}. Consequently, $\delta(G)\geq 5$, and Puleo's Lemma~2.7(a), (b), and~(e)~\cite{Puleo} give the independence of vertices of degree at most six, the complement structure of their links, and the degree-nine lower bound for subsumers of degree-five vertices.

For the link edge count, the complement restriction gives at least nine edges for five-vertex links and at least twelve edges for six-vertex links. Links of order at least eight have at least eight edges by minimum degree two. At order seven, equality with seven edges would mean every component is a cycle; however, a maximum alternating matching on an even cycle leaves a matching, and on an odd cycle leaves a two-edge path and isolated edges, both of which have vertex cover number at most the matching size, contradicting Lemma~\ref{lem:basic}.
\end{proof}

\begin{lemma}\label{lem:vc-and-two-matchings}
Let $v\in V(G)$.
\begin{enumerate}
\item[(a)] For every matching $M$ in $L_v$, the graph $L_v-M$ has vertex cover number at least three, and in particular contains a matching of size two.
\item[(b)] After deleting the union of any two matchings of size two (or a single matching of size two) from $L_v$, the residual graph retains a matching of size two.
\end{enumerate}
\end{lemma}

\begin{proof}
For (a), suppose $Q\subseteq V(L_v)$ covers $L_v-M$ with $|Q|\leq 2$. If $|M|\geq 2$, then $M$ and $Q$ give a whole-link reduction (Lemma~\ref{lem:basic}). If $|M|\leq 1$ and $Q$ already covers all edges of $L_v$, minimum degree two forces $|Q|=2$ and every vertex outside $Q$ adjacent to both vertices of $Q$; since $|V(L_v)|\geq 5$, there is a two-edge matching between $Q$ and $V(L_v)\setminus Q$, reducing with $Q$. Otherwise $M$ consists of a single edge $e$ with both endpoints outside $Q$. Choosing a vertex $w\in V(L_v)\setminus(Q\cup e)$ and pairing an edge from $w$ to $Q$ with $e$ produces a two-edge matching whose deletion leaves all edges covered by $Q$, again a reduction. Since any graph without a two-edge matching is a star or triangle (plus isolates) and has vertex cover number at most two, $L_v-M$ contains a matching of size two.

For (b), let $F$ be the intersection of the deleted edge set with $E(L_v)$, so $|F|\leq 4$ and $F$ is a union of two matchings. Suppose $L_v-F$ has no two-edge matching, so its edges form a star or a triangle. For $|V(L_v)|=5$ or $6$, Lemma~\ref{lem:puleo} gives $|E(L_v)|\geq 9$ or $12$, so $|E(L_v-F)|\geq 5$ or $8$, which is too large for a star or triangle on five or six vertices. For $|V(L_v)|\geq 7$, Lemma~\ref{lem:puleo} gives $|E(L_v-F)|\geq 4$, so $L_v-F$ must be a star with some center $c$. By $\delta(L_v)\geq 2$, each of the at least six vertices in $V(L_v)\setminus\{c\}$ is incident with $F$.

Every component of $F$ is a path or an even cycle. If $F$ contains a three-edge matching $M$, then $L_v-M$ consists of the star at $c$ and at most one remaining edge of $F$, which is covered by $c$ and at most one further vertex, contradicting Lemma~\ref{lem:basic}. Otherwise, since $F$ has at most four edges, is a union of two matchings, and covers at least six vertices, $F$ must consist of two disjoint two-edge paths, and $c$ lies outside both paths (as $|V(L_v)|\geq 7$ and all vertices other than $c$ are covered by $F$, forcing $|V(L_v)|=7$). By $\delta(L_v)\geq 2$, all four path endpoints are adjacent to $c$. Choosing one edge from each path together with an edge from $c$ to an unused path endpoint forms a three-edge matching whose deletion leaves edges covered by $c$ and the two path midpoints, again contradicting Lemma~\ref{lem:basic}.
\end{proof}

By Kandasamy's theorem~\cite{Kandasamy}, $\Delta(G)\geq 11$. Fix a vertex $a\in V(G)$ of maximum degree $d=\Delta(G)$, set $H=G[N(a)]$, and let $X=V(G)\setminus(\{a\}\cup N(a))$. Write $\mu$ for graph matching number. By Lemma~\ref{lem:basic}, $\delta(H)\geq 2$.

\begin{lemma}\label{lem:tb}
If a graph $R$ on $d$ vertices satisfies $\mu(R)\leq m$, then there exists $S\subseteq V(R)$ of size $s=|S|\leq m$ such that $R-S$ has $q\geq s+d-2m$ odd components. If $i$ of these odd components are isolated vertices, then
\[
d-s\geq i+3(q-i), \qquad i\geq 2s+d-3m.
\]
Each odd component of $R-S$ has order at least the smallest positive odd integer at least $\delta(R)-s+1$.
\end{lemma}

\begin{proof}
This follows immediately from the Tutte--Berge formula~\cite{Schrijver}: $q\leq d-s$ forces $s\leq m$, and $d-s\geq i+3(q-i)\geq 3(s+d-2m)-2i$ rearranges to $i\geq 2s+d-3m$.
\end{proof}

\begin{lemma}\label{lem:isolates}
For any $S\subseteq V(H)$, let $i$ be the number of isolated vertices of $H-S$. Then $i=2,3,4,5$ isolated vertices supply respectively $4,6,7,8$ edge-disjoint triangles in $G$, all of whose edges are disjoint from all triangles through $a$ supported on edges of the nonisolated part of $H-S$. In particular, $i\leq 4$.
\end{lemma}

\begin{proof}
The isolated vertices of $H-S$ form an independent set, and all of their neighbors in $G$ lie in $\{a\}\cup S\cup X$. At the first three isolates, choose two-edge link matchings successively avoiding the previously used base edges (using Lemma~\ref{lem:vc-and-two-matchings}(b) at the second and third isolates), yielding six edge-disjoint triangles using six base edges inside $\{a\}\cup S\cup X$. At the fourth and fifth isolates, choose one unused link edge each (possible since every link has at least eight edges by Lemma~\ref{lem:puleo}). This gives $4,6,7,8$ edge-disjoint triangles for $i=2,3,4,5$, respectively. Every edge of these triangles is either a spoke from an isolate to $\{a\}\cup S\cup X$ or a base edge inside $\{a\}\cup S\cup X$, and is therefore disjoint from every triangle through $a$ supported on an edge in the nonisolated part of $H-S$.
\end{proof}

\section{The maximum-degree vertex is universal}

In this section we assume $X\neq\varnothing$ and derive a contradiction.

\begin{lemma}\label{lem:mu5}
We have $\mu(H)=5$.
\end{lemma}

\begin{proof}
Suppose first that $\mu(H)\geq 6$, and fix a six-edge matching $M$ in $H$, giving six triangles through $a$. If $X$ is independent, then for any $x\in X$ we have $L_x\subseteq H$, and $L_x-(M\cap E(L_x))$ has a two-edge matching by Lemma~\ref{lem:vc-and-two-matchings}(a), adding two triangles through $x$. If $X$ contains an edge $xy$, choose a triangle $xyz$ containing $xy$; it uses no edge of $H$ and avoids $a$. By Lemma~\ref{lem:vc-and-two-matchings}(a), $\{y,z\}\cap V(L_x)$ does not cover $L_x-(M\cap E(L_x))$, so there is an edge $uv\in E(L_x)\setminus M$ avoiding $y,z$, and $xuv$ together with $xyz$ adds two triangles. Both cases give $\nu(G)\geq 8$, so $\mu(H)\leq 5$.

Now suppose $\mu(H)\leq 4$, and apply Lemma~\ref{lem:tb} to $H$ with $m=4$, giving $s\leq 4$, $q\geq s+d-8$, and $i\geq 2s+d-12$.
\begin{itemize}
\item For $s\geq 3$, we have $i\geq 5$, contradicting Lemma~\ref{lem:isolates}.
\item For $s=2$, we have $i\geq d-8\geq 3$. If $i\geq 5$, Lemma~\ref{lem:isolates} gives a contradiction. If $i=4$, at least $d-6\geq 5$ nonisolated vertices remain in $H-S$, supplying an apex triangle alongside the seven isolate triangles. If $i=3$, then $d=11$ and $q\geq 5$, leaving at least two nontrivial odd components on the remaining six vertices; taking one edge from each gives two apex triangles alongside the six isolate triangles.
\item For $s=1$, minimum degree two rules out isolates, so $3(d-7)\leq d-1$, impossible for $d\geq 11$.
\item For $s=0$, every component has order at least three ($\delta(H)\geq 2$), so $3(d-8)\leq d$ forces $d\in\{11,12\}$, with component orders $3,3,5$ or $3,3,3,3$, respectively.
\end{itemize}
In the last case, we claim each three-vertex component $uvw$ supplies three edge-disjoint triangles using its three edges and distinct third vertices from $\{a\}\cup X$. Assign to each edge of $uvw$ its list of common neighbors in $\{a\}\cup X$; all three lists contain $a$. If two lists equal $\{a\}$, say for $uv$ and $uw$, then $\{a,v,w\}$ is a triangle component of $L_u$, contradicting Lemma~\ref{lem:basic}. If one list equals $\{a\}$ and the list union has size two, the other two lists contain a common $x\in X$, which is adjacent to all three vertices and belongs to the singleton list as well. If all three lists equal $\{a,x\}$, then $W=\{a,v,w,x\}$, $M=\{av,wx\}$, and $Q=\{w,x\}$ give a scoped reduction at $u$ (only $x$ can have link edges leaving $W$). These cases exhaust all failures of Hall's condition. Applying distinct representatives to two three-vertex components yields six edge-disjoint triangles, and a two-edge matching in the remaining component(s) (order five, or two triangles) adds two apex triangles for eight in total. Thus $\mu(H)=5$.
\end{proof}

\begin{lemma}\label{lem:X-small}
The set $X$ is a clique of size at most two. Consequently, $\delta(H)\geq 3$.
\end{lemma}

\begin{proof}
Fix a five-edge matching $M$ in $H$, giving five apex triangles.

First suppose $x,y\in X$ are nonadjacent. By Lemma~\ref{lem:vc-and-two-matchings}(a), $L_x-(M\cap E(L_x))$ has a two-edge matching $\{e_1,e_2\}$, yielding two triangles through $x$ avoiding $M$. Because $xy\notin E(G)$, these two triangles can share edges with a triangle through $y$ only via $e_1,e_2$. Since $L_y-(M\cap E(L_y))$ has vertex cover number at least three (Lemma~\ref{lem:vc-and-two-matchings}(a)), it has at least three edges, so an edge $e_3\notin\{e_1,e_2\}$ gives a third new triangle through $y$, for eight in total. Hence $X$ is a clique.

If $|X|\geq 6$, then $X$ contains three edge-disjoint triangles, giving eight alongside the five apex triangles.

If $|X|=5$, then by $\delta(G)\geq 5$ there is an edge $xh$ with $x\in X$ and $h\in V(H)$. This edge lies in at least two triangles, at most one of which uses an edge of $M$ incident with $h$. Choose a triangle on $xh$ avoiding $M$; it uses at most one edge inside $X$. Since $K_5$ minus an edge contains two edge-disjoint triangles (for instance $012$ and $034$ when $13$ is removed), we obtain three new triangles avoiding $M$, for eight in total.

If $|X|=4$, choose $x\in X$ and let $T=X\setminus\{x\}$ be a triangle. Since $d_G(x)\geq 5$, the link $L_x$ contains at least two vertices in $H$, and every $h\in N(x)\cap V(H)$ retains positive degree in $R=L_x-(M\cap E(L_x))-E(T)$. We claim $R$ has a two-edge matching. Otherwise, the edges of $R$ form a star or a triangle.
\begin{itemize}
\item If $R$ is a star with center $c\in V(T)$, then $c$ together with one other vertex of $T$ covers $L_x-(M\cap E(L_x))$, contradicting Lemma~\ref{lem:vc-and-two-matchings}(a).
\item If $R$ is a star with center $c\in V(H)$, choose $h\in (N(x)\cap V(H))\setminus\{c\}$. Minimum degree two in $L_x$ forces $h$ to be incident with an edge of $M'=M\cap E(L_x)$, so $|M'|\geq 1$. Adding one edge $e_T\in E(T)$ to $M'$ gives a matching $M''$ in $L_x$ of size $|M'|+1\geq 2$ (as $M'\subseteq E(H)$ and $e_T\subseteq E(T)$ are vertex-disjoint), while $c$ and the vertex of $T$ opposite $e_T$ cover $L_x-M''$, contradicting Lemma~\ref{lem:basic}.
\item If $R$ is a triangle, then every vertex of $N(x)\cap V(H)$ belongs to this triangle, so $|N(x)\cap V(H)|\in\{2,3\}$. Because every pair of vertices in $N(x)\cap V(H)$ is already an edge of $R$, we have $M\cap E(L_x)=\varnothing$. Thus $L_x$ consists either of two disjoint triangles (when $|N(x)\cap V(H)|=3$) or of two triangles sharing a vertex (when $|N(x)\cap V(H)|=2$), both contradicting Lemma~\ref{lem:puleo}.
\end{itemize}
Hence $R$ has a two-edge matching, yielding two triangles through $x$ disjoint from $T$ and $M$, for eight triangles in total.

If $|X|=3$, write $X=\{x,y,z\}$. For each $w\in X$, let $E_w=E(H[N(w)])\setminus M$. Because the other two vertices of $X$ cannot cover $L_w-(M\cap E(L_w))$ (Lemma~\ref{lem:vc-and-two-matchings}(a)), each $E_w$ is nonempty. If $|E_x\cup E_y\cup E_z|\geq 2$, we can select distinct edges $e_1\in E_{w_1}$ and $e_2\in E_{w_2}$ at distinct vertices $w_1,w_2\in X$, giving triangles $w_1+e_1$, $w_2+e_2$, and $xyz$ avoiding $M$, for eight in total. Otherwise $E_x=E_y=E_z=\{uv\}$. Since $d_G(x)\geq 5$, the vertex $x$ sees $u,v$ and at least one further vertex $h\in V(H)\setminus\{u,v\}$. In $L_x-(M\cap E(L_x))$, the vertex $h$ has positive degree but lies on no edge with both endpoints in $H$, so $h$ is adjacent to another vertex of $X$, say $y$. Then the triangles $xuv$, $xyh$, and $yzu$ are edge-disjoint and avoid $M$, again giving eight triangles.

Thus $|X|\leq 2$. If some $v\in V(H)$ had $d_H(v)=2$, then $\delta(G)\geq 5$ would force $|X|=2$ and $v$ adjacent to both vertices of $X$, giving a five-vertex link $L_v$ in which $a$ has two nonneighbors, contradicting Lemma~\ref{lem:puleo}. Therefore $\delta(H)\geq 3$.
\end{proof}

\begin{lemma}\label{lem:exposable-sees-X}
Every vertex $v\in V(H)$ missed by some five-edge matching $M$ of $H$ is adjacent to every $x\in X$.
\end{lemma}

\begin{proof}
Suppose $xv\notin E(G)$ for some $x\in X$ and $v\in V(H)$ missed by a five-edge matching $M$ of $H$. By Lemma~\ref{lem:vc-and-two-matchings}(a), $L_v-(M\cap E(L_v))$ has vertex cover number at least three, so deleting $a$ leaves $L_v-\{a\}-(M\cap E(L_v))$ with vertex cover number at least two, and hence at least two edges. Since $v$ is unmatched by $M$, every edge of $L_v-\{a\}-(M\cap E(L_v))$ forms a triangle through $v$ edge-disjoint from the five apex triangles of $M$.

Meanwhile, $L_x-(M\cap E(L_x))$ contains a two-edge matching $\{e_1,e_2\}$, giving two triangles through $x$ avoiding $M$. Since $xv\notin E(G)$, these two triangles can exhaust the available base edges at $v$ only if $E(L_v-\{a\}-(M\cap E(L_v)))=\{e_1,e_2\}$. In that case, because $|V(L_x)|\geq 5$, there is a vertex $w\in V(L_x)$ outside the four endpoints of $e_1,e_2$, and $w$ has positive degree in $L_x-(M\cap E(L_x))$. Taking an edge incident with $w$ together with whichever of $e_1,e_2$ is disjoint from it gives a new two-edge matching in $L_x-(M\cap E(L_x))$ leaving the other edge free for $v$. In either case we obtain three triangles avoiding the five apex triangles, a contradiction.
\end{proof}

\begin{proposition}\label{prop:universal}
The maximum-degree vertex $a$ is universal: $X=\varnothing$.
\end{proposition}

\begin{proof}
Continue assuming $X\neq\varnothing$, and fix $x\in X$. Apply the Gallai--Edmonds decomposition~\cite{Schrijver} to $H$ (where $\mu(H)=5$ by Lemma~\ref{lem:mu5}): let $D$ be the set of vertices missed by some maximum matching of $H$, and set $S=N_H(D)\setminus D$, with $s=|S|\leq 5$. Then $H[D]$ has $q=s+d-10$ odd factor-critical components, the remaining vertices form even components of $H-S$ with perfect matchings, and a maximum matching inside $H-S$ has $5-s$ edges. By Lemma~\ref{lem:exposable-sees-X}, every vertex in every odd component of $H[D]$ is adjacent to $x$.

Let $i$ be the number of isolated vertices in $H-S$, so $i\geq 2s+d-15$ by Lemma~\ref{lem:tb}. By Lemma~\ref{lem:isolates}, $i\leq 4$. If $i=4$, then $2s+d-15\leq 4$ forces $s\leq 4$ (as $d\geq 11$), so $5-s\geq 1$, and a matching in the nonisolated part of $H-S$ adds at least one apex triangle to the seven isolate triangles. If $i=3$, then $2s+d-15\leq 3$ forces $s\leq 3$, so $5-s\geq 2$, adding at least two apex triangles to the six isolate triangles. Thus $i\leq 2$.

We analyze the remaining values of $s$:
\begin{itemize}
\item \textbf{$s\geq 3$:} From $2\geq i\geq 2s+d-15$ and $d\geq 11$, the only possibility is $s=3$, $d=11$, and $i=2$. Then $q=4$ and the other two odd components of $H[D]$ are triangles. Each triangle gives two edge-disjoint triangles by extending two of its edges through $a$ and $x$, respectively, and the two isolates supply four more triangles (Lemma~\ref{lem:isolates}) using disjoint edges, for eight in total.
\item \textbf{$s=2$:} Since $\delta(H)\geq 3$ (Lemma~\ref{lem:X-small}), $i=0$ and every odd component has order at least three. Thus $3(d-8)\leq d-2$, forcing $d=11$ and three triangular components in $H[D]$. In each triangle, every vertex has a neighbor in $S$ (by $\delta(H)\geq 3$); since $|S|=2$, two vertices share a neighbor in $S$. Extending their edge through that vertex of $S$ and the other two edges through $a$ and $x$ yields three triangles per component, for nine edge-disjoint triangles in total.
\item \textbf{$s=1$:} Here $\delta(H-S)\geq 2$, so $i=0$ and every component of $H-S$ (even or odd) has order at least three (in particular, no even component has order two). From $3(d-9)\leq d-1$ we obtain $d\in\{11,12,13\}$. Each odd component of order three is a triangle complete to the sole vertex of $S$, yielding three triangles using that vertex, $a$, and $x$.

We claim that any odd component $B$ of $H[D]$ of order $b\in\{5,7\}$ supports a near-perfect matching $M_1$ (at $a$) and an edge-disjoint two-edge matching $M_2$ (at $x$), yielding $(b-1)/2+2$ triangles ($4$ for $b=5$, $5$ for $b=7$). Start with a near-perfect matching $M_1$ of $B$; then $B-M_1$ has minimum degree at least one. If $B-M_1$ has no two-edge matching, it must be a spanning star centered at the vertex $c$ missed by $M_1$, with $M_1$ pairing all leaves. Replacing one leaf-pair edge $u_1v_1\in M_1$ with $cu_1$ gives a new near-perfect matching $M_1'$, and taking $M_2'=\{u_1v_1,cu_2\}$ (using a leaf $u_2$ from another pair) gives the desired edge-disjoint two-edge matching.

The possible component-order partitions of $H-S$ are: for $d=11$, $(3,7)$, $(5,5)$, or $(3,3,4)$ (where the even order-$4$ component supplies a two-edge matching at $a$); for $d=12$, $(3,3,5)$; for $d=13$, $(3,3,3,3)$. Every case yields at least eight edge-disjoint triangles.
\item \textbf{$s=0$:} Since $\delta(H)\geq 3$, every component of $H$ has order at least four, and every odd component has order at least five. From $5(d-10)\leq d$ we obtain $d\in\{11,12\}$.

An odd component $B$ of order five has $\delta(B)\geq 3$, so it is $K_5$ minus at most a matching. Labeling its vertices $0,1,2,3,4$ with any missing edges among $\{01,23\}$, we obtain five edge-disjoint triangles: $024$, the matching $\{03,14\}$ extended at $a$, and the matching $\{12,34\}$ extended at $x$.

An odd component $B$ of order $b\in\{7,11\}$ supports a near-perfect matching $M_1$ (at $a$) and an edge-disjoint three-edge matching $M_2$ (at $x$), yielding $6$ triangles for $b=7$ and $8$ for $b=11$. Indeed, let $R=B-M_1$, which satisfies $\delta(R)\geq 2$. If $\mu(R)\leq 2$, applying Lemma~\ref{lem:tb} to $R$ gives $S'\subseteq V(R)$ of size $s'\leq 2$ with $q'\geq s'+b-4$ odd components in $R-S'$. For $s'=0$ or $1$, odd components have order at least three and require more than $b-s'$ vertices. For $s'=2$, all $b-2$ vertices of $R-S'$ are isolated in $R-S'$, hence adjacent in $R$ to both vertices of $S'$. Because $B$ is simple and $\delta(B)\geq 3$, none of these $b-2$ vertices can have an edge of $M_1$ into $S'$, so $M_1$ must pair all $b-2$ isolates with one another, which is impossible since $b-2$ is odd.

The possible component-order partitions of $H$ are: for $d=12$, $(5,7)$ (yielding $5+6=11$ triangles); for $d=11$, $(5,6)$, $(7,4)$, or $(11)$ (where even components of orders $6$ and $4$ supply perfect matchings of sizes $3$ and $2$ at $a$, yielding $5+3=8$, $6+2=8$, and $8$ triangles, respectively).
\end{itemize}
All constructions in different components use disjoint edges, yielding a contradiction in every case. Hence $X=\varnothing$.
\end{proof}

\section{A universal vertex of degree at least thirteen}

By Proposition~\ref{prop:universal}, $a$ is universal, $V(G)=\{a\}\cup V(H)$, and $\delta(H)\geq 4$. Moreover, $H$ is connected: otherwise, the cones over its components form edge-disjoint proper subgraphs partitioning all triangles of $G$, and summing $\tau\leq 2\nu$ over them contradicts the choice of $G$. Since $a$ is universal, $\mu(H)\leq 7$.

\begin{lemma}\label{lem:mu-at-least-6}
If $d\geq 12$, then $\mu(H)\geq 6$.
\end{lemma}

\begin{proof}
Suppose $\mu(H)\leq 5$. Apply Lemma~\ref{lem:tb} to $H$ with $m=5$ and $\delta(H)\geq 4$, giving $s\leq 5$, $q\geq s+d-10$, and $i\geq 2s+d-15$. Connectivity rules out $s=0$ ($q\geq d-10\geq 2$). For $s=1,2,3$, the minimum odd orders $5,3,3$ multiplied by $s+d-10$ exceed $d-s$ for all $d\geq 12$. For $s\geq 4$, we get $i\geq 5$, contradicting Lemma~\ref{lem:isolates}.
\end{proof}

\begin{proposition}\label{prop:deg13plus}
We have $d\leq 12$.
\end{proposition}

\begin{proof}
Suppose $d\geq 13$. By Lemma~\ref{lem:mu-at-least-6}, $\mu(H)\in\{6,7\}$.

\medskip
\noindent\textbf{Case 1: $\mu(H)=6$.}
Every vertex $v\in V(H)$ missed by a six-edge matching $M$ of $H$ has $d_H(v)\geq 6$. Indeed, if $d_H(v)=4$, then $N_H(v)$ induces $K_4$ minus at most a matching (Lemma~\ref{lem:puleo}) and contains at least two edge-disjoint perfect matchings; any edge of one meets both edges of the other, so $M$ avoids at least one of them, yielding two triangles through $v$ in $H$ disjoint from the six apex triangles of $M$. If $d_H(v)=5$, then $N_H(v)$ induces $K_5$ minus at most a matching; deleting $M$ leaves minimum degree at least two and thus a two-edge matching, again giving two triangles through $v$ avoiding $M$.

Apply the Gallai--Edmonds decomposition to $H$: let $D$ be the set of vertices missed by some maximum matching and $S=N_H(D)\setminus D$, with $s=|S|\leq 6$. Then $H[D]$ has $q=s+d-12$ odd components, every vertex of $D$ has degree at least six in $H$, and all neighbors of $D$ lie in $D\cup S$. For $s=1,2,3,4,5$, the minimum odd orders $7,5,5,3,3$ multiplied by $s+d-12$ exceed $d-s$ at $d=13$ and for all larger $d$. For $s=6$, all $d-6\geq 7$ remaining vertices are isolated, contradicting Lemma~\ref{lem:isolates}. Thus $s=0$. Since $D\neq\varnothing$ and $H$ is connected, $D=V(H)$ and $q=1$, forcing $d=13$ and $\delta(H)\geq 6$.

If $\nu(H)\leq 1$, minimality of $G$ gives $\tau(H)\leq 2$; together with twelve spokes at $a$ (which hit every triangle through $a$), this gives a cover of $G$ of size at most fourteen. Otherwise, choose two edge-disjoint triangles $T_1,T_2$ in $H$ and set $R=H-E(T_1)-E(T_2)$. Every vertex outside $T_1\cup T_2$ retains degree at least six in $R$, every vertex in one of $T_1,T_2$ has degree at least four, and at most one shared vertex has degree below four (at least two).

If $\mu(R)\geq 6$, then six apex triangles from $R$ together with $T_1,T_2$ give eight triangles. Suppose $\mu(R)\leq 5$. Lemma~\ref{lem:tb} gives $s\leq 5$, $q\geq s+3$, and $i\geq 2s-2$.
\begin{itemize}
\item For $s=0,1$, there are no isolates, and no order-three component can exist (it would require three vertices of degree at most three in $R$); thus odd orders are at least five, requiring at least $15$ or $20$ vertices.
\item For $s=2,3$, at most one vertex (a shared vertex) can be isolated, giving at least $1+3(q-1)>13-s$ vertices.
\item For $s=5$, $i\geq 8$ exceeds the at most six vertices of $T_1\cup T_2$.
\item For $s=4$, we need $i\geq 6$, forcing $T_1$ and $T_2$ to be vertex-disjoint with all six of their vertices isolated in $R-S$. By $\delta(H)\geq 6$, each of these six vertices is adjacent to all four vertices of $S$. The remaining three vertices outside $S$ are unaffected, retain degree at least six in $R$, and have no edges to the six isolates, so they form a third triangle $T_3$ complete to $S$. Extending the three edges of each of $T_1,T_2,T_3$ using three distinct vertices of $S$ yields nine edge-disjoint triangles.
\end{itemize}
This excludes $\mu(H)=6$.

\medskip
\noindent\textbf{Case 2: $\mu(H)=7$.}
Fix a seven-edge matching $M$ in $H$; then $H-M$ is triangle-free. No vertex of $H$ has degree four or five: either degree supplies a five-vertex subset of $H$ inducing $K_5$ minus a matching (Lemma~\ref{lem:puleo}), and after deleting $M$ its complement is a union of two disjoint matchings and hence bipartite, yielding an independent triple in the complement and a triangle in $H-M$. Thus $\delta(H)\geq 6$.

Applying Lemma~\ref{lem:tb} to $H$ with $m=7$ and $\delta(H)\geq 6$ excludes $d\geq 16$: connectivity rules out $s=0$, $s\geq 6$ forces $i\geq 7$ (contradicting Lemma~\ref{lem:isolates}), and for $s=1,2,3,4,5$ the minimum odd orders $7,5,5,3,3$ multiplied by $s+d-14$ exceed $d-s$ for all $d\geq 16$. Hence $d\in\{14,15\}$.

Suppose $H$ contains a triangle $T$, and set $R=H-E(T)$. Only the three vertices of $T$ can have degree below six in $R$, and $\delta(R)\geq 4$. If $\mu(R)\geq 7$, then $T$ and seven apex triangles give eight triangles. Suppose $\mu(R)\leq 6$; Lemma~\ref{lem:tb} gives $s\leq 6$, $q\geq s+d-12$, and $i\geq 2s+d-18$.
\begin{itemize}
\item For $s=0$, no odd component has order three (as $T$ is independent in $R$ and $\delta(R)\geq 4$) or order five (which would contain an unaffected vertex of degree at least six in $R$). Thus every odd component has order at least seven. This excludes $d=15$ ($q\geq 3$). For $d=14$, $R$ has two order-seven components, and $T$ crosses between them because $H$ is connected. At least five vertices in each component lie outside $T$; their degree at least six makes each complete to its component, giving a $K_5$ in each component. Each cone over $K_5$ is $K_6$, which contains an octahedron $K_{2,2,2}$ packing four edge-disjoint triangles (choosing one vertex from each of three pairs with even parity sum). Together the two cones pack eight edge-disjoint triangles.
\item For $s=1$ (minimum odd order $5$) and $s=3$ (minimum odd order $3$), the vertex lower bounds exceed $d-s$.
\item For $s=2$ and $d=15$, minimum odd order $3$ requires $15>13$ vertices. For $s=2$ and $d=14$, fitting $q\geq 4$ odd components in $12$ vertices requires four order-three components, whereas every unaffected vertex forces its component to have order at least five, and only three vertices lie in $T$.
\item For $s=4,5$, the lower bound $i\geq 2s+d-18\geq 4$ exceeds the three vertices of $T$.
\item For $s=6$, all $d-6$ vertices of $R-S$ are isolated in $R-S$. At least $(d-6)-3=d-9\geq 5$ of them lie outside $T$ and are therefore isolated in $H-S$ as well, contradicting Lemma~\ref{lem:isolates}.
\end{itemize}
Consequently $H$ is triangle-free. Every triangle of $G$ then uses two spokes at $a$, so any $d-1\leq 14$ spokes cover all triangles of $G$, contradicting $\tau(G)>14$.
\end{proof}

\section{A universal vertex of degree twelve}

\begin{proposition}\label{prop:deg12}
We have $d=11$.
\end{proposition}

\begin{proof}
By Proposition~\ref{prop:deg13plus}, $d\in\{11,12\}$. Suppose $d=12$. By Lemma~\ref{lem:mu-at-least-6}, $\mu(H)=6$, so $H$ has a perfect matching.

We first prove $\delta(H)\geq 6$. Fix a perfect matching $M$ of $H$. If some vertex $v\in V(H)$ has $d_H(v)=5$, then by Lemma~\ref{lem:puleo} its closed neighborhood $N_H[v]$ induces $K_6$ minus a matching. Deleting $M$ leaves a six-vertex graph whose complement is a union of two disjoint matchings—a union of paths and even cycles. The two color classes of each component differ in size by at most one, and since the total order is even, the components combine into a bipartition with both parts of size three. Those two parts yield two vertex-disjoint triangles in $H-M$, giving eight triangles alongside the six apex triangles of $M$.

Similarly, if $d_H(v)=4$ and $N_H(v)$ induces $K_4$, then $N_H[v]$ induces $K_5$. Deleting $M$ leaves a vertex of $N_H[v]$ missed by $M\cap E(H[N_H[v]])$ and a two-edge matching on the remaining four vertices, yielding two triangles in $H-M$ and eight in total.

Suppose now that $d_H(v)=4$ and $N_H(v)=\{p,q,x,y\}$ induces $K_4-pq$ (the only other option under Lemma~\ref{lem:puleo}). Then $x$ and $y$ subsume $v$ in $G$, so by Lemma~\ref{lem:puleo} they have degree at least nine in $G$, hence at least eight in $H$. Also $p$ and $q$ are adjacent to the degree-five vertex $v$, so by Lemma~\ref{lem:puleo} they have degree at least seven in $G$, hence at least six in $H$. Let $O=V(H)\setminus\{v,p,q,x,y\}$, so $|O|=7$. Each of $x,y$ has at least four neighbors in $O$, and each of $p,q$ has at least three neighbors in $O$.

Delete the five edges among $\{p,q,x,y\}$ from $H$ to obtain $R$. Observe that $\delta(R)\geq 4$: $v$ has degree $4$, $p,q$ have degree at least $1+3=4$, $x,y$ have degree at least $1+4=5$, and vertices in $O$ retain degree at least $4$. If $R$ has a perfect matching $M_R$, then $M_R$ uses only one edge inside $N_H[v]$ (incident with $v$). Thus $H-M_R$ induces $K_5$ minus at most two edges on $N_H[v]$, which contains two edge-disjoint triangles: for two disjoint missing edges, extend a matching on four vertices through the fifth; for adjacent missing edges $01,02$, take $034$ and $123$. Together with the six apex triangles of $M_R$, this gives eight triangles.

Hence $R$ has no perfect matching. By Tutte's theorem (using parity since $d=12$ is even), there exists $S\subseteq V(R)$ with $s=|S|\leq 5$ such that $R-S$ has $q\geq s+2$ odd components and $i\geq 2s-3$ isolated vertices.
\begin{itemize}
\item For $s=0$, $R$ is connected because each deleted edge among $\{p,q,x,y\}$ has its endpoints joined through $v$.
\item For $s=1,2,3$, $\delta(R)\geq 4$ gives minimum odd orders $5,3,3$, requiring at least $15,12,15$ vertices, exceeding $12-s$ in each case.
\item For $s=4$, we need $i\geq 5$. If $v\in S$, then $|S\setminus\{v\}|=3$; no vertex of $O$ can be isolated (none sees $v$ and each has degree at least four), nor can $x,y$ (each has at least four neighbors in $O$), leaving at most $p,q$ as isolates. If $v\notin S$, none of $p,q,x,y$ can be isolated, so at least four vertices of $O$ are isolated in $R-S$ (and hence in $H-S$) and complete to $S$. Taking four such isolates $I_0$, setting $W=I_0\cup S$, matching $I_0$ to $S$, and setting $Q=S$ gives a scoped link reduction at $a$ (Lemma~\ref{lem:basic}).
\item For $s=5$, all seven remaining vertices of $R-S$ are isolated. If $v\notin S$, then $\{p,q,x,y\}\subseteq S$, leaving a set $I\subseteq O$ of six isolates in $H-S$, each with at least four neighbors in $S$. Let $Q=N_H(I)\subseteq S$. Every singleton in $Q$ has a neighbor in $I$, and every subset of $Q$ of size at least two meets every vertex of $I$ (as each isolate misses at most one vertex of $S$), so $|N_I(A)|=6\geq |A|$. Hall's theorem gives a matching $M_0$ saturating $Q$ into $I$, and taking $W=Q\cup I$ with $M_0$ and $Q$ yields a scoped link reduction at $a$.
\item If $s=5$ and $v\in S$, let $S'=S\setminus\{v\}$ ($|S'|=4$) and $I=O\setminus S$. Every vertex of $I$ has degree at least four in $H$ and does not see $v$, so it is complete to $S'$. If $|I|\geq 4$, matching four vertices of $I$ to $S'$ gives a scoped link reduction at $a$ as above. Otherwise $|I|=3$, so $S'\subseteq O$ and $p,q,x,y\notin S$. Since $x,y$ are isolated in $R-S$ and each has at least four neighbors in $O$, both $x$ and $y$ are complete to $S'$. Now take four edge-disjoint triangles in the octahedron on $\{a,v,p,q,x,y\}$ omitting $av,xy,pq$, add a triangle $xyz$ with $z\in S'$, and add three apex triangles from a matching between $I$ and $S'$. All eight triangles are edge-disjoint.
\end{itemize}
This eliminates degree-four vertices and proves $\delta(H)\geq 6$.

Now if $\nu(H)\leq 1$, minimality gives $\tau(H)\leq 2$, and adding eleven spokes at $a$ covers all triangles of $G$ with at most thirteen edges. Otherwise, choose two edge-disjoint triangles $T_1,T_2$ in $H$ and let $R=H-E(T_1)-E(T_2)$. Only a possible shared vertex of $T_1,T_2$ has degree below four in $R$ (at least two), and every unaffected vertex retains degree at least six. If $R$ has a perfect matching, it yields six apex triangles alongside $T_1,T_2$.

Suppose $R$ has no perfect matching, giving $s\leq 5$, $q\geq s+2$, and $i\geq 2s-3$.
\begin{itemize}
\item For $s=0$, every odd component has order at least five. An order-five component $C$ is also impossible: by $\delta(H)\geq 6$, all five vertices of $C$ must lose edges across the cut $(C,V(H)\setminus C)$, whereas each triangle $T_1,T_2$ can meet at most two vertices of $C$ with crossing edges. Thus every odd component has order at least seven, contradicting $q\geq 2$.
\item For $s=1$, no isolates or order-three components are possible (at most one vertex has degree below four in $R$), so $5\cdot 3=15>11$.
\item For $s=2$, at most one isolate is possible, so the only fit in ten vertices is one isolate and three order-three components. Every vertex of an order-three component has degree at most four in $R$ and must be affected, requiring ten affected vertices where at most five exist.
\item For $s=3$, at most one isolate is possible, giving $1+3(4)=13>9$.
\item For $s=5$, all seven remaining vertices must be isolated, exceeding the at most six affected vertices.
\item For $s=4$, we need $i\geq 5$, and every isolate is an affected vertex. If $T_1,T_2$ are disjoint, the isolates contain a full triangle $uvw$ and a disjoint edge $xy$ of the other triangle, with all five vertices adjacent to all of $S$ (by $\delta(H)\geq 6$). If $T_1,T_2$ share a vertex $v$, all five affected vertices are isolated; write $uvw$ for one triangle and $vxy$ for the other, so $u,w,x,y$ see all of $S$ and $v$ sees at least two vertices of $S$. In either case, let $S=\{s_1,s_2,s_3,s_4\}$ with $v$ adjacent to $s_1,s_2$. Extend $vu,vw,uw,xy$ through $s_1,s_2,s_3,s_4$, respectively, and add the four apex triangles on $s_1w,s_2x,s_3y,s_4u$. All eight triangles are edge-disjoint.
\end{itemize}
This contradiction excludes $d=12$ and proves $d=11$.
\end{proof}

\section{A universal vertex of degree eleven}

By Propositions~\ref{prop:universal} and~\ref{prop:deg12}, $a$ is universal and $H$ is a connected graph on $d=11$ vertices with $\delta(H)\geq 4$. We eliminate $\delta(H)=4$, $\delta(H)=5$, and $\delta(H)\geq 6$ in turn.

\begin{lemma}\label{lem:d11-deg4}
We have $\delta(H)\geq 5$.
\end{lemma}

\begin{proof}
Suppose $v\in V(H)$ has $d_H(v)=4$ with neighbors $p,q,x,y$, and let $O=V(H)\setminus\{v,p,q,x,y\}$, so $|O|=6$. By Lemma~\ref{lem:puleo}, $H[\{p,q,x,y\}]$ is either $K_4$ or $K_4-pq$. Each vertex universal on $\{p,q,x,y\}$ subsumes $v$ in $G$, so by Lemma~\ref{lem:puleo} it has degree at least eight in $H$ and at least four neighbors in $O$. In particular, $C=N_O(x)\cap N_O(y)$ has size at least two. The octahedron on $\{a,v,p,q,x,y\}$ omitting $av,xy,pq$ packs four edge-disjoint triangles avoiding these three edges.

\medskip
\noindent\textbf{Case 1: $H[\{p,q,x,y\}]=K_4$.}
All four of $p,q,x,y$ subsume $v$, so both $\{x,y\}$ and $\{p,q\}$ have common neighbors in $O$ (possibly the same vertex). Extending $xy$ and $pq$ through these common neighbors adds two triangles disjoint from the octahedron and from the cone $a+O$, giving six triangles and forcing $\mu(H[O])\leq 1$.

Cover the four spokes from $a$ to $\{p,q,x,y\}$ and the six edges among $\{p,q,x,y\}$, at cost ten. Every remaining triangle uses an edge of $H[O]$. If $|E(H[O])|\leq 4$, adding all edges of $H[O]$ gives a cover of size at most fourteen. Otherwise, since $\mu(H[O])\leq 1$ on six vertices, $H[O]$ is a five-edge star with some center $w\in O$. Every remaining triangle then uses an edge from $w$ to $\{a,p,q,x,y\}$. If at most two of $\{p,q,x,y\}$ see $w$, adding those at most three spokes to $w$ gives a cover of size at most thirteen. Otherwise, choose three neighbors of $w$ in $\{p,q,x,y\}$ together with $a$. Each chosen neighbor of $v$ sees at least three of the five leaves of the star, and $a$ sees all five. Hall's condition assigns four distinct leaves to these four vertices, forming four edge-disjoint triangles with $w$ that are edge-disjoint from the four octahedron triangles, for eight in total.

\medskip
\noindent\textbf{Case 2: $H[\{p,q,x,y\}]=K_4-pq$.}
The octahedron together with $xyz$ (for any $z\in C$) packs five triangles edge-disjoint from the cone $a+O$, so $\nu(a+O)\leq 2$. Start a cover with the four spokes from $a$ to $\{p,q,x,y\}$ and the five edges among $\{p,q,x,y\}$, at cost nine. Every remaining triangle either lies in $a+O$ or extends an edge of $F=H[O]$ through one of $p,q,x,y$. If $|E(F)|\leq 5$, adding all edges of $F$ gives a cover of size at most fourteen. Hence $|E(F)|\geq 6$, which forces $\mu(F)=2$ (since $\mu(F)\leq\nu(a+O)\leq 2$ and $\mu(F)\leq 1$ allows at most five edges on six vertices).

Call an edge of $F$ \emph{compulsory} if it lies in every two-edge matching of $F$. Any edge $e\in E(F)$ extended to a triangle through $p$ or $q$ must be compulsory; otherwise the five triangles (octahedron plus $xyz$), the extension of $e$, and two apex triangles on a two-edge matching of $F$ avoiding $e$ give eight triangles. The same argument applies to an extension of $e$ through $x$ or $y$ whenever $C$ contains a vertex $z\notin e$. Since $|C|\geq 2$, at most one edge of $F$ can contain $C$ (call it \emph{exceptional}). Moreover, $F$ has at most one compulsory edge: if a two-edge matching were unique, every other edge of $F$ would have to join its two edges, and at most two such cross-edges could exist without creating a second two-edge matching, contradicting $|E(F)|\geq 6$.

If $F$ is triangle-free, then $F$ is bipartite: an odd cycle on six vertices would be a chordless $C_5$, and any edge from the sixth vertex to the $C_5$ would yield a three-edge matching, contradicting $|E(F)|\geq 6$. By K\H{o}nig's theorem, $F$ has a vertex cover of size two. Adding the two spokes from $a$ to this vertex cover and the at most two (compulsory or exceptional) edges of $F$ extended through $p,q,x,y$ yields a cover of size at most $9+2+2=13$.

If $F$ contains a triangle $T$, then $\mu(F-E(T))\leq 1$ (else $T$ plus two apex triangles gives $\nu(a+O)\geq 3$). Since $|E(F-E(T))|\geq 3$, the graph $F-E(T)$ is either a triangle $T'$ or a star with center $c$ and at least three leaves.
\begin{itemize}
\item If $F-E(T)$ is a triangle $T'$, then $T$ and $T'$ share at most one vertex, and there is no compulsory edge: to avoid any edge $e$, choose a different edge in its triangle and a disjoint edge in the other triangle (picking the edge opposite the shared vertex on one side when needed).
\item If $F-E(T)$ is a star with center $c\notin V(T)$, there is again no compulsory edge: to avoid a star edge, take another star edge and an edge of $T$ avoiding its leaf; to avoid an edge of $T$, take a star edge to a leaf outside $T$ if one exists, or otherwise note that $V(T)\cup\{c\}$ induces $K_4$, which has three edge-disjoint perfect matchings.
\end{itemize}
In both subcases, $a+O$ has a cover of size at most $2\nu(a+O)\leq 4$ by minimality, and at most one exceptional edge covers all extensions through $p,q,x,y$, giving a total cover of size at most $9+4+1=14$. Finally, if $F-E(T)$ is a star with center $c\in V(T)$, then the edge $f\in E(T)$ opposite $c$ is compulsory and every other edge of $F$ is incident with $c$. Adding $f$, the spoke $ac$, and the at most one exceptional edge gives a cover of size at most $9+3=12$. This rules out $d_H(v)=4$.
\end{proof}

\begin{lemma}\label{lem:d11-deg5}
We have $\delta(H)\geq 6$.
\end{lemma}

\begin{proof}
Suppose $v\in V(H)$ has $d_H(v)=5$, and let $O=V(H)\setminus N_H[v]$, so $|O|=5$. Because $a$ is universal in the six-vertex link $L_v$, Lemma~\ref{lem:puleo} implies that the complement of $H[N_H(v)]$ has maximum degree at most one and edge count different from two; on five vertices, this forces $H[N_H(v)]$ to be $K_5$ or $K_5-pq$.

If $H[N_H(v)]=K_5$, then $\{a,v\}\cup N_H(v)$ induces $K_7$, whose edges partition into seven triangles via the Fano plane. Any edge inside $O$ adds an eighth triangle through $a$. If $O$ is independent, then by $\delta(H)\geq 5$ (Lemma~\ref{lem:d11-deg4}) every vertex of $O$ is complete to $N_H(v)$. Taking $W=O\cup N_H(v)$, a five-edge matching between $O$ and $N_H(v)$, and $Q=N_H(v)$ gives a scoped link reduction at $a$ (Lemma~\ref{lem:basic}).

If $H[N_H(v)]=K_5-pq$, label the Fano decomposition of $K_7$ on $\{a,v\}\cup N_H(v)$ so that $apq$ is one of the seven triples. Omitting $apq$ leaves six edge-disjoint triangles and leaves the spokes $ap$ and $aq$ unused. By $\delta(H)\geq 5$, each of $p,q$ has at least one neighbor in $O$. If $p,q$ have distinct neighbors $z_1,z_2\in O$, then $apz_1$ and $aqz_2$ add two triangles for eight in total. Otherwise both $p$ and $q$ have the same unique neighbor $z\in O$. Any edge $xy$ inside $O\setminus\{z\}$ then adds two triangles $apz$ and $axy$. If $O\setminus\{z\}$ is independent, every vertex of $O\setminus\{z\}$ can only see $z$ and the three vertices of $N_H(v)\setminus\{p,q\}$, giving degree at most four in $H$ and contradicting Lemma~\ref{lem:d11-deg4}.
\end{proof}

\begin{proof}[Proof of Theorem~\ref{thm:main}]
By Propositions~\ref{prop:universal} and~\ref{prop:deg12} and Lemmas~\ref{lem:d11-deg4} and~\ref{lem:d11-deg5}, $a$ is universal and $H$ has $d=11$ vertices with $\delta(H)\geq 6$.

If $\nu(H)\leq 3$, minimality of $G$ gives a triangle cover $C_H$ of $H$ with $|C_H|\leq 6$. The eleven-vertex triangle-free graph $H-C_H$ has an independent set $I$ of size three: a vertex of degree at least three has three independent neighbors, while maximum degree at most two yields a greedy independent set of size at least $\lceil 11/3\rceil=4$. Adding the eight spokes from $a$ to $V(H)\setminus I$ to $C_H$ covers all triangles of $G$ with at most fourteen edges, a contradiction.

Thus $\nu(H)\geq 4$. Choose three edge-disjoint triangles $T_1,T_2,T_3$ in $H$, and let $R=H-E(T_1)-E(T_2)-E(T_3)$. If $\mu(R)\geq 5$, then five apex triangles from $R$ together with $T_1,T_2,T_3$ yield eight triangles. Suppose $\mu(R)\leq 4$; Lemma~\ref{lem:tb} gives $S\subseteq V(R)$ with $s=|S|\leq 4$ and $q\geq s+3$ odd components in $R-S$.

Call a vertex \emph{shared} if it lies in at least two of $T_1,T_2,T_3$. Since the three triangles are edge-disjoint, there are at most three shared vertices, and only shared vertices can have degree below four in $R$. If a vertex lies in all three triangles, it is the only shared vertex. If there are three shared vertices, each is shared by a distinct pair of triangles, so every pair of shared vertices forms a deleted triangle edge and the three shared vertices are independent in $R$.

We examine each $s\in\{0,1,2,3\}$:
\begin{itemize}
\item \textbf{$s=0,1$:} Only a vertex shared by all three triangles can be isolated, so $i\leq 1$. Any order-three odd component would have to consist of three shared vertices, impossible since three shared vertices are independent in $R$. At $s=1$, at least three odd components of order at least five cannot fit in ten vertices. At $s=0$, at least $11-9=2$ (indeed $11-7=4$ when $i=1$) vertices are unaffected by the deleted edges and retain degree at least six in $R$; the only count fitting in eleven vertices is one isolate plus two order-five components, neither of which can contain an unaffected vertex.
\item \textbf{$s=2$:} Only shared vertices can be isolated, so $i\leq 3$. Fitting $q\geq 5$ odd components in nine vertices forces three isolates and two order-three components. Since all three shared vertices exist, $|V(T_1\cup T_2\cup T_3)|=6$, leaving five unaffected vertices in $H$ and hence at least three in $R-S$. Their degree at least four in $R-S$ prevents them from lying in either order-three component.
\item \textbf{$s=3$:} At most three isolates exist, so $q\geq 6$ odd components require at least $3+3(3)=12>8$ vertices.
\end{itemize}

Finally, consider $s=4$. All seven vertices of $U=V(H)\setminus S$ are isolated in $R-S$, so $E(H[U])\subseteq E(T_1)\cup E(T_2)\cup E(T_3)$. Because $\delta(H)\geq 6$, we have $\delta(H[U])\geq 2$ and $|E(H[U])|\geq 7$. If any $T_j$ contained a vertex of $S$, then $|E(H[U])|\leq 3+3+1=7$, forcing $H[U]$ to be a $2$-regular graph on seven vertices containing two triangles, which is impossible. Thus all three triangles lie inside $U$, and $H[U]=T_1\cup T_2\cup T_3$ has seven vertices and nine edges with minimum degree two. Letting $n_r$ count vertices in $r$ triangles gives $n_1+n_2+n_3=7$ and $n_1+2n_2+3n_3=9$, so $n_2+2n_3=2$. Hence $H[U]$ is either a chain of three triangles (the middle triangle sharing distinct vertices with the two end triangles) or three triangles sharing a common center. Every ordinary vertex (degree two in $H[U]$) sees all four vertices of $S$, and every shared vertex of degree four in $H[U]$ sees at least two vertices of $S$.

\begin{itemize}
\item \textbf{Chain of three triangles:} Omit the edge joining the two shared vertices, and extend the remaining eight edges of the chain through $\{a\}\cup S$. Each shared vertex has at least three available vertices in $\{a\}\cup S$, and every ordinary vertex has all five. First extend the two edges from the shared vertices to their common ordinary neighbor using two distinct available vertices. Next, at each shared vertex independently, extend its other two incident chain edges using two distinct remaining available vertices. Finally, extend each of the two edges between ordinary vertices using any of the five vertices in $\{a\}\cup S$ avoiding the two vertices already used at its endpoints. All eight triangles are edge-disjoint.
\item \textbf{Three triangles with common center $v$:} Write one triangle as $vxy$. If $v$ has no neighbor in $S$, then $L_v$ in $G$ consists of a star centered at $a$ together with three disjoint leaf-pair edges (the three edges opposite $v$). Taking these three edges as $M$ and $Q=\{a\}$ gives a whole-link reduction at $v$ (Lemma~\ref{lem:basic}). Otherwise $v$ has a neighbor $s\in S$. Replace $vxy$ by $vxs$ and $vya$, keep the two other triangles at $v$, and extend the four ordinary vertices of those two triangles with $a$ and the four distinct vertices of $S$ (matching one vertex of $S$ to each of the four ordinary vertices). This yields eight edge-disjoint triangles.
\end{itemize}
Every case leads to a contradiction, completing the proof.
\end{proof}

\begin{thebibliography}{9}
\bibitem{Kandasamy}
A.~Kandasamy, \emph{Tuza's conjecture for graphs of maximum degree at most 10},
\href{https://arxiv.org/abs/2610.15269}{arXiv:2610.15269}, 2026.

\bibitem{Puleo}
G.~J.~Puleo, \emph{Tuza's conjecture for graphs with maximum average degree less than 7},
\href{https://doi.org/10.1016/j.ejc.2015.03.006}{European J. Combin. 49 (2015), 134--152}.

\bibitem{Schrijver}
A.~Schrijver, \emph{Combinatorial Optimization: Polyhedra and Efficiency},
\href{https://www.lamsade.dauphine.fr/~cornaz/Enseignement/M2_MODO/DATA/A3.PDF}{Section~24}, Springer, 2003.

\bibitem{PackingSix}
\emph{Tuza's conjecture for triangle packing number at most six}, companion manuscript,
\href{run:packing-six.pdf}{packing-six.pdf}.
\end{thebibliography}

\end{document}
